\documentclass[border=10pt]{standalone}
\usepackage{fontspec}
\setmainfont{Noto Serif}
\newfontfamily\bengalifont{Noto Serif Bengali}[Script=Bengali]
\usepackage[Latin,Bengali]{ucharclasses}
\setTransitionsFor{Bengali}{\begingroup\bengalifont}{\endgroup}
\usepackage{newunicodechar}
\newunicodechar{।}{{\bengalifont\char"0964}}
\newunicodechar{॥}{{\bengalifont\char"0965}}

\usepackage[table]{xcolor}
\usepackage{booktabs,tabularx,array}
\definecolor{bapilrule}{HTML}{1F2937}
\definecolor{bapilstripe}{HTML}{F3F4F6}
\renewcommand{\arraystretch}{1.28}
\setlength{\tabcolsep}{5pt}

\begin{document}
\begin{minipage}{22cm}
{\large\bfseries সারণি ২৩।}\ {\large পাইলট কর্পাসের জন্য জনসমক্ষে প্রকাশিত ডেটা-উৎসের প্রাথমিক তালিকা}\par\vspace{7pt}
\small
\rowcolors{2}{bapilstripe}{white}
\begin{tabularx}{\linewidth}{@{}>{\hsize=1.189\hsize\raggedright\arraybackslash}X>{\hsize=1.189\hsize\raggedright\arraybackslash}X>{\hsize=1.245\hsize\raggedright\arraybackslash}X>{\hsize=0.623\hsize\raggedright\arraybackslash}X>{\hsize=0.453\hsize\raggedright\arraybackslash}X>{\hsize=1.302\hsize\raggedright\arraybackslash}X@{}}
\arrayrulecolor{bapilrule}\toprule
\rowcolor{white}\textbf{উৎস} & \textbf{রক্ষক} & \textbf{বিষয়বস্তু} & \textbf{ফরম্যাট} & \textbf{শ্রেণি} & \textbf{প্রবেশাধিকার-সংক্রান্ত টীকা} \\
\midrule
বাংলাদেশের আইন (bdlaws.minlaw.gov.bd) & লেজিসলেটিভ ও সংসদ বিষয়ক বিভাগ & বাংলা ও ইংরেজিতে আইন ও অধ্যাদেশ & HTML & P0 & পুনর্ব্যবহারের শর্ত নিশ্চিত করা; একসঙ্গে রপ্তানির অনুরোধ (সিদ্ধান্ত D3) \\
বাংলাদেশ গেজেট (dpp.gov.bd) & মুদ্রণ ও প্রকাশনা অধিদপ্তর & গেজেট প্রজ্ঞাপন & PDF, প্রায়ই স্ক্যানকৃত & P0 & OCR-এর মান একটি বড় ঝুঁকি (R3) \\
জাতীয় বাজেট নথি (mof.gov.bd) & অর্থ বিভাগ & বাজেট বক্তৃতা, মধ্যমেয়াদি বাজেট কাঠামো, বরাদ্দ & PDF, স্প্রেডশিট & P0 & পাইলট P2-এর ভিত্তি \\
উন্নয়ন পরিকল্পনা (plancomm.gov.bd) & পরিকল্পনা কমিশন & পঞ্চবার্ষিক পরিকল্পনা, প্রেক্ষিত পরিকল্পনা, এডিপি & PDF, স্প্রেডশিট & P0 & পরিকল্পনা সংশোধনের মধ্যে সংস্করণ ব্যবস্থাপনা \\
সরকারি পরিসংখ্যান (bbs.gov.bd) & বাংলাদেশ পরিসংখ্যান ব্যুরো & বর্ষগ্রন্থ ও জরিপ প্রতিবেদন & PDF, স্প্রেডশিট & P0 & শুধু সামষ্টিক তথ্য; মাইক্রোডেটা বর্জিত \\
উন্মুক্ত সরকারি ডেটা (data.gov.bd) & আইসিটি বিভাগ ও এটুআই & যন্ত্র-পাঠযোগ্য ডেটাসেট & CSV, স্প্রেডশিট, API & P0 & ডেটাসেটভিত্তিক লাইসেন্স যাচাই \\
জাতীয় তথ্য বাতায়ন (bangladesh.gov.bd) & আইসিটি বিভাগ ও এটুআই & দপ্তরের তথ্য, সেবা, বিজ্ঞপ্তি & HTML & P0 & বিপুল পরিমাণ; মন্ত্রণালয়ভিত্তিক অগ্রাধিকার \\
জাতীয় সংসদ (parliament.gov.bd) & বাংলাদেশ জাতীয় সংসদ সচিবালয় & বিল, বিতর্ক, কমিটির প্রতিবেদন & PDF & P0 & স্ক্যানকৃত ও ডিজিটালের মিশ্রণ \\
ই-নথি রেকর্ড & এটুআই ও মন্ত্রণালয়সমূহ & প্রশাসনিক নথি & অভ্যন্তরীণ ব্যবস্থা & P1 থেকে P3 & পাইলটের আওতাবহির্ভূত \\
BAPIL প্রকাশনা (এই রিপোজিটরি) & BAPIL & হোয়াইট পেপার, নির্বাহী সারসংক্ষেপ, ব্লুপ্রিন্ট, গবেষণা & Markdown, PDF & P0 & রেফারেন্স বাস্তবায়নের প্রদর্শনী কর্পাস \\
\bottomrule
\end{tabularx}
\par\vspace{5pt}{\footnotesize\itshape গ্রহণের আগে প্রতিটি সংস্থার সঙ্গে রক্ষক ও প্রবেশাধিকারের শর্ত নিশ্চিত করতে হবে। কোনো উৎস তালিকাভুক্ত করা মানে তার রক্ষকের সঙ্গে কোনো চুক্তি বোঝায় না।}
\par\vspace{3pt}{\scriptsize\color{gray} উৎস: BAPIL কারিগরি ব্লুপ্রিন্ট, সংস্করণ ১.০ (২০২৬)। CC BY 4.0।}
\end{minipage}
\end{document}
